\documentclass[11pt]{article}
\usepackage[a4paper,margin=18mm]{geometry}
\usepackage{fontspec}
\setmainfont{Latin Modern Roman}
\usepackage{amsmath,amssymb,amsthm,amscd,graphicx,color}
\usepackage{xurl}
\usepackage[unicode,hidelinks]{hyperref}
\allowdisplaybreaks
\setlength\emergencystretch{3em}
\usepackage{mathtools} % AMS Math fonts and symbols
\usepackage{slashed} % Slashed symbols e.g. to typeset Dirac operators
\usepackage{mathrsfs}

\usepackage{tikz}
\usetikzlibrary{cd} % Commutative diagrams

% Flattened unchanged original author macros from rudismacros.sty
\makeatletter
\RequirePackage{xparse}

\NewDocumentCommand{\blank}{}{{-}}
\NewDocumentCommand{\parensup}{m}{\textup{(}#1\textup{)}}


% Standard number sets
\NewDocumentCommand{\numbers}{m}{\mathbb{#1}}
\NewDocumentCommand{\C}{}{\numbers{C}}
\NewDocumentCommand{\Z}{}{\numbers{Z}}
\NewDocumentCommand{\N}{}{\numbers{N}}
\NewDocumentCommand{\Q}{}{\numbers{Q}}
\NewDocumentCommand{\R}{}{\numbers{R}}

% Eulers number and imaginary unit
\NewDocumentCommand{\eu}{}{\mathrm{e}}
\NewDocumentCommand{\iu}{}{\mathrm{i}}


% Support
\DeclareMathOperator{\supp}{supp}

% Differential e.g. to typeset dx in integrals
\NewDocumentCommand{\D}{}{\mathop{}\!\mathrm{d}}

% Differential geometry
\DeclareMathOperator{\scal}{scal}

% Classical Groups
\DeclareMathOperator{\GL}{GL}
\DeclareMathOperator{\SL}{SL}
\DeclareMathOperator{\Orth}{O}
\DeclareMathOperator{\SO}{SO}
\DeclareMathOperator{\U}{U}
\DeclareMathOperator{\SU}{SU}

% Homology and K-theory
\DeclareMathOperator{\HZ}{H}
\DeclareMathOperator{\K}{K}
\DeclareMathOperator{\KO}{KO}
\NewDocumentCommand{\cpt}{}{\mathrm{c}}
\DeclareMathOperator{\ph}{ph}

% A Hat
\NewDocumentCommand{\Ahat}{}{\hat{\mathrm{A}}}
\NewDocumentCommand{\AhatDeg}{}{\Ahat\!-\!\deg}
\NewDocumentCommand{\AhatClass}{}{\hat{\mathbf{A}}}

% Spheres, disks, balls
\NewDocumentCommand{\Sphere}{}{\mathrm{S}}
\NewDocumentCommand{\Torus}{}{\mathrm{T}}
\NewDocumentCommand{\Disk}{}{\mathrm{D}}
\NewDocumentCommand{\Ball}{}{\mathrm{B}}
\NewDocumentCommand{\Nbh}{}{\mathcal{U}}

% Distance and width
\DeclareMathOperator{\width}{width}
\DeclareMathOperator{\dist}{dist}

% Boundary, interior, double
\DeclareMathOperator{\double}{\mathsf{d}\!}
\DeclareMathOperator{\bd}{\partial}
\NewDocumentCommand{\interior}{m}{#1^{\circ}}

% Classifying spaces
\NewDocumentCommand{\Bfree}{}{\mathrm{B}}
\NewDocumentCommand{\Efree}{}{\mathrm{E}}
\NewDocumentCommand{\Eub}{}{\underline{\mathrm{E}}}

% Index theory
\DeclareMathOperator{\ind}{index}
\DeclareMathOperator{\indc}{index_{\mathrm{c}}}
\DeclareMathOperator{\indrel}{indrel}
\NewDocumentCommand{\Dirac}{}{\slashed{\mathfrak{D}}}
\NewDocumentCommand{\SpinBdl}{}{\mathfrak{S}}
\DeclareMathOperator{\clm}{c}
\NewDocumentCommand{\tensgr}{}{\mathbin{\widehat{\otimes}}}

% Cstar-algebras
\NewDocumentCommand{\Cstar}{}{\mathrm{C}^\ast}
\NewDocumentCommand{\red}{}{\mathrm{red}}
\NewDocumentCommand{\CstarRed}{}{\mathrm{C}^\ast_{\mathrm{red}}}
\NewDocumentCommand{\CstarMax}{}{\mathrm{C}^\ast_{\mathrm{max}}}
\NewDocumentCommand{\textCstar}{}{\ensuremath{\Cstar\!}}
\DeclareMathOperator{\Ltwo}{L^2}
\DeclareMathOperator{\ELL}{L}
\NewDocumentCommand{\Cl}{}{\mathrm{Cl}}
\NewDocumentCommand{\mathcstar}{m}{#1}
\DeclareMathOperator{\Ct}{C}
\NewDocumentCommand{\Cb}{}{\Ct_{\mathrm{b}}}
%\DeclareMathOperator{\Cb}{\Ct_{\mathrm{b}}}
\DeclareMathOperator{\Cc}{\Ct_{\mathrm{c}}}
\DeclareMathOperator{\Cz}{\Ct_{0}}
\DeclareMathOperator{\Cgr}{\mathcal{S}}
\DeclareMathOperator{\Lin}{\mathcal{L}}
\DeclareMathOperator{\Kom}{\mathcal{K}}

% Roe Algebras
\NewDocumentCommand \RoeSymbol {o} {
	\mathrm{C}^{\ast}
	\IfNoValueF{#1}{_{#1}}
}
\NewDocumentCommand \Roe {o} {\RoeSymbol[#1]}


% PSC
\DeclareOption{psc}{
  \NewDocumentCommand \StolzBordism {o} {
     \Omega^{
      \mathrm{spin}
      \IfValueT{#1}{,#1}
     }
  }

  \NewDocumentCommand \StolzRel {o} {
     \operatorname{R}^{
      \mathrm{spin}
      \IfValueT{#1}{,#1}
     }
  }

  \NewDocumentCommand \StolzPos {o} {
     \operatorname{Pos}^{
      \mathrm{spin}
      \IfValueT{#1}{,#1}
     }
  }

  \NewDocumentCommand \StolzPartialPSC {o} {
     \operatorname{R}\Omega^{
      \mathrm{spin}
      \IfValueT{#1}{,#1}
     }
  }

  \DeclareMathOperator{\Riem}{\mathcal{R}}
  \DeclareMathOperator{\RiemP}{\mathcal{R}^{+}}

}

\DeclareOption{higsonroe}{
  \DeclareMathOperator{\StructGp}{S}
}
\ProcessOptions\relax

\makeatother


\numberwithin{equation}{section}

\newtheorem{Theorem}{Theorem}[section]
\newtheorem{Corollary}[Theorem]{Corollary}
\newtheorem{Lemma}[Theorem]{Lemma}
\newtheorem{Proposition}[Theorem]{Proposition}
\newtheorem{Conjecture}[Theorem]{Conjecture}
 { \theoremstyle{definition}
\newtheorem{Definition}[Theorem]{Definition}
\newtheorem{Note}[Theorem]{Note}
\newtheorem{Example}[Theorem]{Example}
\newtheorem{Remark}[Theorem]{Remark} }

\newcommand*{\Mishchenko}{Mi\v{s}\v{c}enko}

\begin{document}
\begin{center}\Large A closed two-sided pi1-injective hypersurface with nonzero Rosenberg index in a complete spin n-manifold forces infimum scalar curvature <=-n(n-1)tan\textasciicircum{}2(nd/4) when its width-d neighbourhood has scalar curvature >=n(n-1) and 0<d<2pi/n\end{center}
\noindent\textbf{Stable ID:} 1652\par
\noindent\textbf{Authors:} Rudolf Zeidler\par
\noindent\textbf{Original work:} Width, Largeness and Index Theory\par
\noindent\textbf{Version:} Official SIGMA 16 (2020) article 127, DOI 10.3842/SIGMA.2020.127; journal PDF and native TeX acquired 2026-09-30, exact bytes pinned by SHA256\par
\noindent\textbf{Selected task scope:} Exactly Theorem1.4 restricted to n>=2: X is spin without boundary, g complete, M is a closed codimension-one submanifold with trivial normal bundle and pi1(M)->pi1(X) injective; alpha(M) in KO\_(n-1)(C*pi1M) is nonzero. If scal\_g>=n(n-1) throughout the neighbourhood \{dist(p,M)<=d/2\} with 0<d<2*pi/n, then inf\_X scal\_g<=-n(n-1)*tan(nd/4)\textasciicircum{}2. The proof uses the explicitly flat Mishchenko bundle on the covering specified by the author. Separate width/iso-enlargeability corollaries and arbitrary nonflat coefficient bundles are excluded.\par
\noindent\textbf{Difficulty:} H2: The obligation is a quantitative negative scalar-curvature bound away from a positively curved band, rather than only a qualitative obstruction to uniformly positive curvature. The author designs a smooth global Callias potential by a truncated Riccati equation, controls it inside and outside the band and forces a uniform operator-square lower bound that contradicts the hypersurface index. This combines a sharp geometric inequality with Hilbert-C*-module index theory beyond standard qualifying Dirac-operator material.\par
\noindent\textbf{Editorial boundary note (not author text):} The full original Theorem 1.4 is retained; the selected n>=2 flat Mishchenko-bundle hypersurface application is specified on the cover. The complete new Callias-potential argument of Theorem 3.1 and the actual covering-space/signed-distance reduction are included, together with the Rosenberg-index and flat-bundle definitions. The Riccati cutoff, comparison, exterior lower bound, Friedrich/Lichnerowicz inequality and index contradiction are all author text. Prior Callias index identification remains cited standard background. Separate width/largeness consequences are excluded.\par
\noindent\textbf{Source:} \url{https://sigma-journal.com/2020/127/}\quad DOI: \url{https://doi.org/10.3842/SIGMA.2020.127}\par
\noindent\textbf{License and attribution:} Copyright retained by the named authors. Published by SIGMA under CC BY 4.0: \url{https://creativecommons.org/licenses/by/4.0/}. Source: \url{https://sigma-journal.com/about.html}. Adaptation consists of standalone excerpt formatting, cover and numbering restoration; original author language and mathematical text are retained.\par
\noindent\textbf{Editorial symbol notice (not author text):} Editorial scope notice: Section3’s general setup mentions a metric connection on E, but the selected Theorem1.4 application explicitly uses the flat Mishchenko bundle, defined as flat in Section2. The scalar-only Dirac estimate is being used solely in that author specified flat application. The source wording remains unchanged and no assertion for arbitrary nonflat E is included.\par
\medskip\hrule\medskip\noindent\textbf{Author text begins below.}\par
\setcounter{section}{1}
\setcounter{subsection}{0}
\setcounter{subsubsection}{0}
\setcounter{Theorem}{3}
\setcounter{equation}{0}
% Original source lines 109--123
This was prompted by observing that both cases are based on the same index-theoretic technique, namely an index theorem for certain \emph{Callias-type} operators, suggesting that there should be a more general geometric statement underlying both.
This is achieved in Theorem~\ref{thm:quantitative-codim1-obstr} below.
We show that in the presence of a~hypersurface \(M \subset X\) with an index-theoretic obstruction, positive scalar curvature \(\geq \sigma > 0\) in a~neighborhood \(\{p \in X \mid d(p, M) \leq d/2 \}\) of width \(d\) forces the scalar curvature somewhere else to be negative in a~quantifiable way depending on \(d\) and \(\sigma\).
We formulate this for the normalized case \(\sigma = n(n-1) = \scal_{\Sphere^n}\).

\begin{Theorem} \label{thm:quantitative-codim1-obstr}
 Let \(X\) be a spin \(n\)-manifold without boundary and \(M \subseteq X\) a closed codimension one submanifold with trivial normal bundle such that the induced map \(\pi_1 M \to \pi_1 X\) is injective.
 Suppose that the Rosenberg index \(\alpha(M) \in \KO_{n-1}(\Cstar \pi_1 M)\) does not vanish.

 Let \(g\) be a complete Riemannian metric on \(X\) which has scalar curvature bounded below by \(n(n-1)\) in a neighborhood of \(M\) of width \(d < 2\pi/n\).
 Then
 \begin{equation} \label{eq:negativity}
 \inf_{p \in X} \scal_g(p) \leq - n(n-1) \tan\left(\frac{nd}{4}\right)^2.
 \end{equation}
\end{Theorem}

\setcounter{section}{1}
\setcounter{subsection}{0}
\setcounter{subsubsection}{0}
\setcounter{Theorem}{6}
\setcounter{equation}{1}
% Original source lines 151--167
\section{The Rosenberg index} \label{sec:Rosenberg-index}
In this section, we provide a brief survey of the aspects of the Rosenberg index relevant to our statements and open problems.
More material on these matters can be found in the original articles of Rosenberg~\cite{Rosenberg:PSCNovikovI,Rosenberg:PSCNovikovII,Rosenberg:PSCNovikovIII} and in \cite{RosenbergStolz:PSCSurgeryConnections,Schick:ICM,Stolz:ManifoldsOfPSC,Stolz98Concordance}.

The starting point is that the spinor Dirac operator on a spin manifold is closely related to the scalar curvature via the Schr\"odinger--Lichnerowicz formula \(\Dirac^2 = \nabla^\ast \nabla + \frac{\scal}{4}\).
Together with the Atiyah--Singer index theorem, this was first used by Lichnerowicz~\cite{Lichnerwociz:Spineurs} to show that the \(\Ahat\)-genus is an obstruction to the existence of a positive scalar curvature metric on a closed spin manifold \(M\).
Later, this connection between the Dirac operator and scalar curvature was exploited further using variations of the index invariant.
For closed spin manifolds, the most sophisticated such invariant is the \emph{Rosenberg index} \(\alpha(M) \in \KO_n(\Cstar \pi_1 M)\) which takes values in the (real) K-theory of the group \textCstar-algebra \(\Cstar \pi_1 M\).
Notable examples of manifolds with non-vanishing Rosenberg index (but which in general do not have non-vanishing \(\Ahat\)-genus) include (area-)enlargeable spin manifolds~\cite{Hanke-Kotschick-Roe-Schick:Coarse-Top-Enlargeability, HankeSchick:Enlargeable,HankeSchick:EnlargeableInfinite}.

We will now describe two (equivalent) points of view towards the Rosenberg index.
The first is as the index of the spinor Dirac operator of \(M\) twisted by the ``universal flat bundle'' on~\(M\).
The latter is the \emph{Mi\v{s}\v{c}enko bundle} \(\mathcal{L}_M \to M\) which is the flat bundle of finitely generated projective Hilbert-\(\Cstar \Gamma\)-modules associated to the representation of \(\Gamma \coloneqq \pi_1 M\) on \(\Cstar \Gamma\) by left-multiplication.
Index theory extends to the situation of elliptic operators with coefficients in \textCstar-algebras yielding index invariants in corresponding operator K-theory groups.
This goes back to Mi\v{s}\v{c}enko and Fomenko~\cite{Mishchenko-Fomenko:Index}; for more recent expositions of the technical background relevant to our situation we refer to~\cite{Ebert:EllipticRegularityDirac,HankePapeSchick:CodimensionTwoIndex}.
The Rosenberg index \(\alpha(M) \in \KO_n(\Cstar \pi_1 M)\) is then simply the (K-theoretic) index of the twisted Dirac operator on \(\SpinBdl_M \otimes \mathcal{L}_M\) using the flat connection induced on the \Mishchenko\ bundle \(\mathcal{L}_M\), and \(\SpinBdl_M\) is the spinor bundle.
The Schr\"odinger--Lichnerowicz formula extends to this situation so that the non-vanishing of \(\alpha(M) \in \KO_n(\Cstar \Gamma)\) is still an obstruction to the existence of a positive scalar curvature metric on \(M\).

\setcounter{section}{2}
\setcounter{subsection}{0}
\setcounter{subsubsection}{0}
\setcounter{Theorem}{4}
\setcounter{equation}{1}
% Original source lines 249--340
\section{Callias-type operators and codimension one obstructions} \label{sec:codim-1-obstr}
In this section, we will provide a proof of Theorem~\ref{thm:quantitative-codim1-obstr}.
This is a simplified synthesis of the previous approaches from \cite{Cecchini:CalliasTypePSC}, \cite[Section~6]{Cecchini:LongNeck} and \cite[Section~2]{Zeidler:band-width-estimates}.
We rely on the same index-theoretic setup which has its roots in work of Higson~\cite{Higson:CobordismInvariance} and Bunke~\cite{Bunke:RelativeIndexCallias}.

In this section, we fix an arbitrary unital Real \textCstar-algebra \(A\).
For the purposes of the main theorem of this note, one can keep the case \(A = \Cstar \Gamma\) with \(\Gamma\) some discrete group in mind.

Let \(W\) be a complete spin \(n\)-manifold together with a proper smooth \(1\)-Lipschitz function \(x \colon W \to \R\).
Fix some regular value \(a \in \R\) of \(x\) and set \(M \coloneqq x^{-1}(a)\).
Suppose furthermore that \(W\) is endowed with a smooth bundle \(E \to W\) of finitely generated projective Hilbert-\(A\)-modules furnished with a metric connection.
Let \(f \colon \R \to \R\) be some non-decreasing smooth function satisfying \(f(0) = 0\) and such that \(f^2 - f'\) is bounded below by a positive constant outside a~compact subset of~\(\R\).
In particular, the latter property includes the following two special cases:
\begin{enumerate}\itemsep=0pt
 \item[(i)] the function \(f\) is proper and Lipschitz,
 \item[(ii)] or there exists \(R > 0\) such that \(f(x) = \lambda_-\) for all \(x \leq -R\) and \(f(x) = \lambda_+\) for all \(x \geq R\), where \(\lambda_- < 0 < \lambda_+\) are constants.
\end{enumerate}
Denote the \(\Cl_n\)-linear spinor bundle of \(W\) by \(\SpinBdl_W\) and let \(\Dirac_{W,E}\) be the Dirac operator twisted by \(E\).
We then consider the operator
\[
 B = \Dirac_{W, E} \tensgr 1 + f(x) \tensgr \epsilon
\]
acting on compactly supported smooth sections of \(\SpinBdl_W \tensgr E \tensgr \Cl_{0,1}\).
Then \(B\) is an essentially self-adjoint regular Fredholm operator.
It has an index \(\ind(B) \in \KO_{n-1}(A)\) which satisfies
\[
 \ind(B) = \ind(\Dirac_{M, E|_M}) \in \KO_{n-1}(A).
\]
Moreover, we have the formula
\begin{equation}
B^2 = \Dirac^2_{W,E} \tensgr 1 + f'(x) \clm(\D x) \tensgr \epsilon + f(x)^2 \geq \Dirac^2_{W,E} \tensgr 1 + f(x)^2 - f'(x),
\label{eq:callias-estimate}
\end{equation}
where \(\clm\) denotes the Clifford action and we have used \(\|\mathrm{d} x\| \leq 1\) and \(f' \geq 0\).
For all of this we refer, e.g., to \cite[Section~6]{Cecchini:LongNeck}, \cite[Appendix~A]{Zeidler:band-width-estimates}.

\begin{Theorem}\label{thm:quantitative-codim-1}
 In the above setup, suppose that \(\ind(\Dirac_{M, E|_M}) \neq 0 \in \KO_{n-1}(A)\).
 Let \(I \subseteq \R\) be an interval of length \(d < 2 \pi / n\) such that the scalar curvature of \(W\) is bounded below by \(n(n-1)\) on the subset \(x^{-1}(I)\).
 Then
 \begin{equation*}
 \inf_{p \in W} \scal(p) \leq - n(n-1) \tan\left(\frac{nd}{4}\right)^2.
 \end{equation*}
 \end{Theorem}
 \begin{proof}
 We can assume without loss of generality that \(I = [-d/2,d/2]\).
 Then suppose, by contraposition, that there exists \(\sigma < \tan\big(\frac{nd}{4}\big)^2\) such that \(\scal(p) \geq -n(n-1) \sigma\) for all \(p \in W\).
 Choose \(0 < \tilde{d} < d\) such that \(\sigma < \tan\big(\frac{n\tilde{d}}{4}\big)^2 < \tan\big(\frac{nd}{4}\big)^2\) and \(0 < r < n/2\) such that
 \begin{equation}
 \frac{n^2 \sigma}{4} < r^2 \tan\big(r\tilde{d}/2\big)^2. \label{eq:choose-r}
 \end{equation}
 Choose a smooth function \(\xi \colon \R^2 \to \R\) which satisfies \(0 \leq \xi(x,y) \leq y^2 + r^2\) for all \((x,y) \in \R^2\), \(\xi(x,y) = y^2 + r^2\) for \(|x| \leq \tilde{d}/2\) and \(\xi(x,y) = 0\) for \(|x| \geq d/2\).
 Let \(f\) be the unique maximal solution of the initial value problem
 \begin{gather*}
 f'(x) = \xi(x,f(x)), \\
 f(0) = 0.
 \end{gather*}
 By construction, \(f\) is a smooth non-decreasing function defined on some interval containing \(0\).
 We first claim that \(f\) is defined on the entire real line.
 To this end, consider the auxilliary function \(\bar{f} \colon (-\pi/(2 r), \pi/(2r)) \to \R\), \(\bar{f}(x) = r \tan(r x)\).
 Note that this is the unique solution of \(\bar{f}'(x)= \bar{f}(x)^2 + r^2\), \(\bar{f}(0) = 0\) and that \(\pi/(2r) > \pi /n > d/2\).
 By the choice of \(\xi\), \(f'(x) \leq f(x)^2 + r^2\) which implies that \(|f| \leq |\bar{f}|\) as long as both are defined.
 This implies that \(f\) exists at least for all \(|x| \leq d/2\).
 But \(\xi(x,y) = 0\) for \(|x| \geq d/2\), so \(f\) must continue for all times and remain constant on each component of \(\R \setminus (-d/2,d/2)\).
 Moreover, by monotonicity and because \(f(x) = \bar{f}(x)\) for \(|x| \leq \tilde{d}/2\), we obtain the estimate
 \begin{equation}
 |f(x)| \geq |f(\pm \tilde{d}/2)| = \bar{f}(\tilde{d}/2) = r \tan\big(r \tilde{d} / 2\big) \qquad \text{for all \(|x| \geq d/2\).}
 \label{eq:outside-bound}
 \end{equation}

 Now we consider the operator
 \(
 B = \Dirac \tensgr 1 + f(x) \tensgr \epsilon
 \)
 as above, where \(\Dirac = \Dirac_{W,E}\).
 Then \linebreak using~\eqref{eq:callias-estimate} together with the Schr\"odinger--Lichnerowicz formula and the Friedrich estimate (compare~\cite[Section~3.2]{Cecchini:LongNeck}), we obtain
 \begin{align*}
 B^2 &\geq \Dirac^2 \otimes 1 + f(x)^2 - f'(x) \\
 &\geq {\frac{n \scal}{4(n-1)} + f(x)^2 - f'(x)}\ \eqqcolon \Psi.
 \end{align*}
 We denote the latter function by \(\Psi \colon W \to \R\) and will show that it is bounded below by a positive constant to finish the proof.
 By construction of \(f\), we always have \(f(x)^2 - f'(x) \geq -r^2\).
 For \(|x| \leq d/2\), this implies together with the scalar curvature bound on \(x^{-1}(I)\) that
 \[
 \Psi(p) \geq \frac{n^2}{4} - r^2 > 0 \qquad \text{if \(|x(p)| \leq d/2\).}
 \]
 Moreover, for \(|x| \geq d/2\), we have \(f'(x) = 0\) and so \eqref{eq:choose-r} and \eqref{eq:outside-bound} imply
 \begin{align*}
\Psi(p) \geq - \frac{n^2 \sigma}{4} + f(x(p))^2 \geq - \frac{n^2 \sigma}{4} + r^2 \tan\big(r \tilde{d} / 2\big)^2 > 0 \qquad \text{if \(|x(p)| \geq d/2\).}
 \end{align*}
 This proves that \(B^2 \geq \inf\limits_{p \in W} \Psi(p) > 0\) and hence \(\ind(\Dirac_{M, E|_M}) = \ind(B) = 0\).
 \end{proof}

\setcounter{section}{3}
\setcounter{subsection}{0}
\setcounter{subsubsection}{0}
\setcounter{Theorem}{2}
\setcounter{equation}{3}
% Original source lines 351--351
 Note that this implies Theorem~\ref{thm:quantitative-codim1-obstr} from the introduction by taking \(W\) to be the connected covering of \(X\) such that \(\pi_1 W = \pi_1 M\), the bundle \(E = \mathcal{L}_W\) to be the \Mishchenko-bundle, and the function \(x\) to be (a smooth approximation of) the signed distance function to \(M \subseteq W\).
\begin{thebibliography}{99}
\bibitem[3]{Bunke:RelativeIndexCallias}
Bunke U., A {$K$}-theoretic relative index theorem and {C}allias-type {D}irac
 operators, \href{https://doi.org/10.1007/BF01460989}{\textit{Math. Ann.}} \textbf{303} (1995), 241--279.

\bibitem[4]{Cecchini:CalliasTypePSC}
Cecchini S., Callias-type operators in {$C^*$}-algebras and positive scalar
 curvature on noncompact manifolds, \href{https://doi.org/10.1142/S1793525319500687}{\textit{J.~Topol. Anal.}} \textbf{12}
 (2020), 897--939, \href{https://arxiv.org/abs/1611.01800}{arXiv:1611.01800}.

\bibitem[5]{Cecchini:LongNeck}
Cecchini S., A long neck principle for Riemannian spin manifolds with positive
 scalar curvature, \href{https://doi.org/10.1007/s00039-020-00545-1}{\textit{Geom. Funct. Anal.}} \textbf{30} (2020), 1183--1223,
 \href{https://arxiv.org/abs/2002.07131}{arXiv:2002.07131}.

\bibitem[9]{Ebert:EllipticRegularityDirac}
Ebert J., Elliptic regularity for {D}irac operators on families of noncompact
 manifolds, \href{https://arxiv.org/abs/1608.01699}{arXiv:1608.01699}.

\bibitem[16]{Hanke-Kotschick-Roe-Schick:Coarse-Top-Enlargeability}
Hanke B., Kotschick D., Roe J., Schick T., Coarse topology, enlargeability, and
 essentialness, \href{https://doi.org/10.24033/asens.2073}{\textit{Ann. Sci. \'Ec. Norm. Sup\'er.~(4)}} \textbf{41}
 (2008), 471--493, \href{https://arxiv.org/abs/0707.1999}{arXiv:0707.1999}.

\bibitem[17]{HankePapeSchick:CodimensionTwoIndex}
Hanke B., Pape D., Schick T., Codimension two index obstructions to positive
 scalar curvature, \href{https://doi.org/10.5802/aif.3000}{\textit{Ann. Inst. Fourier (Grenoble)}} \textbf{65} (2015),
 2681--2710, \href{https://arxiv.org/abs/1402.4094}{arXiv:1402.4094}.

\bibitem[18]{HankeSchick:Enlargeable}
Hanke B., Schick T., Enlargeability and index theory, \href{https://doi.org/10.4310/jdg/1175266206}{\textit{J.~Differential
 Geom.}} \textbf{74} (2006), 293--320, \href{https://arxiv.org/abs/math.GT/0403257}{arXiv:math.GT/0403257}.

\bibitem[19]{HankeSchick:EnlargeableInfinite}
Hanke B., Schick T., Enlargeability and index theory: infinite covers,
 \href{https://doi.org/10.1007/s10977-007-9004-3}{\textit{$K$-Theory}} \textbf{38} (2007), 23--33, \href{https://arxiv.org/abs/math.GT/0604540}{arXiv:math.GT/0604540}.

\bibitem[20]{Higson:CobordismInvariance}
Higson N., A note on the cobordism invariance of the index, \href{https://doi.org/10.1016/0040-9383(91)90024-X}{\textit{Topology}}
 \textbf{30} (1991), 439--443.

\bibitem[26]{Lichnerwociz:Spineurs}
Lichnerowicz A., Spineurs harmoniques, \textit{C.~R.~Acad. Sci. Paris}
 \textbf{257} (1963), 7--9.

\bibitem[27]{Mishchenko-Fomenko:Index}
Mi\v{s}\v{c}enko A.S., Fomenko A.T., The index of elliptic operators over
 {$C^{\ast} $}-algebras, \href{https://doi.org/10.1070/IM1980v015n01ABEH001207}{\textit{Math. USSR Izv.}} \textbf{15} (1980), 87--112.

\bibitem[30]{Rosenberg:PSCNovikovI}
Rosenberg J., {$C^{\ast} $}-algebras, positive scalar curvature, and the
 {N}ovikov conjecture, \href{https://doi.org/10.1007/BF02953775}{\textit{Inst. Hautes \'Etudes Sci. Publ. Math.}}
 (1983), 197--212.

\bibitem[31]{Rosenberg:PSCNovikovII}
Rosenberg J., {$C^\ast$}-algebras, positive scalar curvature and the {N}ovikov
 conjecture.~{II}, in Geometric Methods in Operator Algebras ({K}yoto, 1983),
 \textit{Pitman Res. Notes Math. Ser.}, Vol.~123, Longman Sci. Tech., Harlow,
 1986, 341--374.

\bibitem[32]{Rosenberg:PSCNovikovIII}
Rosenberg J., {$C^\ast$}-algebras, positive scalar curvature, and the {N}ovikov
 conjecture.~{III}, \href{https://doi.org/10.1016/0040-9383(86)90047-9}{\textit{Topology}} \textbf{25} (1986), 319--336.

\bibitem[35]{RosenbergStolz:PSCSurgeryConnections}
Rosenberg J., Stolz S., Metrics of positive scalar curvature and connections
 with surgery, in Surveys on Surgery Theory, {V}ol.~2, \textit{Ann. of Math.
 Stud.}, Vol.~149, Princeton University Press, Princeton, NJ, 2001, 353--386.

\bibitem[38]{Schick:ICM}
Schick T., The topology of positive scalar curvature, in Proceedings of the
 {I}nternational {C}ongress of {M}a\-the\-maticians~-- {S}eoul 2014, {V}ol.~{II},
 Kyung Moon Sa, Seoul, 2014, 1285--1307, \href{https://arxiv.org/abs/1405.4220}{arXiv:1405.4220}.

\bibitem[41]{Stolz98Concordance}
Stolz S., Concordance classes of positive scalar curvature metrics, {P}reprint,
 1998, available at \url{http://www3.nd.edu/~stolz/concordance.ps}.

\bibitem[43]{Stolz:ManifoldsOfPSC}
Stolz S., Manifolds of positive scalar curvature, in Topology of
 High-Dimensional Manifolds ({T}rieste, 2001), \textit{ICTP Lect. Notes},
 Vol.~9, Abdus Salam Int. Cent. Theoret. Phys., Trieste, 2002, 661--709.

\bibitem[46]{Zeidler:band-width-estimates}
Zeidler R., Band width estimates via the {D}irac operator,
 \textit{J.~Differential Geom.}, {t}o appear, \href{https://arxiv.org/abs/1905.08520}{arXiv:1905.08520}.

\end{thebibliography}
\end{document}
